%%%%%%%%%%%%%%%%%%%%%%%%%%%%%%%%%%%%%%%%%%%%%%%%%%%%%%%%%%%%%%%%%%%%%%%%%%%%%%%%
%\section{Introduction}\label{sec:intro}
%%%%%%%%%%%%%%%%%%%%%%%%%%%%%%%%%%%%%%%%%%%%%%%%%%%%%%%%%%%%%%%%%%%%%%%%%%%%%%%%

Parton distribution functions (PDFs) provide a universal description of the proton's constituents (quarks, antiquarks and gluons).
They are critical inputs~\cite{Alekhin:2012ig,Gao:2013xoa,Radescu:2010zz,CooperSarkar:2011aa,Martin:2009iq,Ball:2012cx}
for the discovery of the Higgs boson, the last particle of the Standard
Model, found at the Large Hadron Collider (LHC)
through proton-proton collisions~\cite{CMS:2012nga,ATLAS:2012oga}. Despite 
this great victory, the LHC has many tasks remaining,
and the focus of the future Runs~2--5 will be to search for physics beyond the
Standard Model. In order 
to discriminate new-physics signatures from the Standard-Model background, we need to improve the precision of the latter. Unfortunately, our knowledge of
many Higgs-production cross sections remains dominated by PDF uncertainties. 
Improvement on current PDF uncertainties is important to assist LHC new-physics searches.

In addition to their applications to the energy frontier, PDFs also reveal
nontrivial structure inside the nucleon, such as the momentum and spin distributions 
of partons. Many ongoing and planned experiments at facilities around 
the world, such as Brookhaven and Jefferson Laboratory in the United States, GSI in Germany, J-PARC in
Japan, or a future electron-ion collider (EIC), are set to explore the less-known 
nucleon structures and more. 
In order to distinguish the flavor content of the PDFs, one would need to 
use nuclear data, such as neutrino scattering off heavy nuclei. However, the current understanding
of medium corrections in these cases is limited. Thus, the uncertainty in the strange
PDFs remains large. In some cases, the assumption $\overline{s}(x)=s(x)$
made in global analyses can agree with data due to the large uncertainty. 
At the LHC, strangeness can be extracted through the $W+c$ associated-production channel,
but their results are not yet well-determined. %rather puzzling. 
For example, ATLAS gets
$(s+\overline{s})/(2\overline{d})=0.96^{+0.26}_{-0.30}$ at $Q^2=1.9\mbox{ GeV}^2$ and 
$x=0.23$~\cite{Aad:2014xca}. 
%, where where $q$ or  $\overline{q}$ means the integral of the $q$-flavor quark or antiquark over the total Bjorken-$x$ range $[0,1]$. 
CMS performs a global analysis with deep-inelastic scattering (DIS) data and the muon-charge asymmetry in 
$W$ production at the LHC to extract the ratios 
of the total integral of strange and anti-strange to the sum of the anti-up and -down, 
%(s+\overline{s})/(\overline{u}+\overline{d})=0.52^{+0.18}_{-0.15}$ at $Q^2=20\mbox{ GeV}^2$
finding it to be $0.52^{+0.18}_{-0.15}$ at $Q^2=20\mbox{ GeV}^2$
~\cite{Chatrchyan:2013mza}.
Future high-luminosity studies may help to improve our knowledge of the strangeness. 
In the polarized case, SU(3)-flavor symmetry is often assumed due to lack of precision 
experimental data. We learn from the unpolarized case that this assumption introduces an underestimated uncertainty. 
In addition, there have been long debates concerning how big the intrinsic charm contribution is 
or whether other heavy flavors contribute. Again, the data is too inconclusive to narrow down or discriminate between the various 
proposed QCD models.

Theoretical determination of the parton distributions is complementary to
the experimental effort, especially for those not yet accessible
kinematically in experiments. On the other hand, in order to be useful for
experiments, theoretical calculations need to demonstrate that those already
measured parton distributions can be reproduced within the same approach.
This turns out to be a challenging task. It is rooted to the
nonperturbative nature of parton interactions inside the nucleon. One
hint of this nonperturbative nature can be seen in the parton distributions
extracted from experimental data. Although the net quark number of the nucleon is 3,
the quark and antiquark numbers are both infinite. This implies that there is no
hierarchy in quark-antiquark pair production through gluon emission.
The production of $N+1$ quark-antiquark pairs is as important as the
production of $N$ pairs, so truncation at a finite $N$ is impossible. This
makes the proton effectively an infinite-body system.

Lattice QCD deals with this infinite-body problem by reducing the
continuous spacetime to a discrete lattice, rendering the number of integrals in the partition function finite.
The lattice is defined in a Euclidean spacetime so that Monte Carlo
algorithms can be used to compute these integrals efficiently. The parton distributions are related to
nucleon matrix elements of quark correlators defined on a lightcone in
Minkowski space. Those lightcone correlators become local operators in
Euclidean space and lead to unphysical results (see more detailed discussion
in the next section). 
In principle, this problem can be avoided by working with moments of parton
distributions, which correspond to matrix elements of local operators,
provided all the moments can be computed to recover the whole PDF. In
practice, one can only obtain the first few (about 3) moments due to
operator mixing with lower-dimension operators with coefficients proportional
to inverse powers of the lattice spacing.
In theory, one can design more complicated operators to 
subtract the power divergence arising from the mixing of high-moment operators 
to get to even higher moments. 
However, the operator renormalization gets significantly more complicated and 
the correlators suffer from signal-to-noise problems as well. 
In recent years more and more lattice-QCD nucleon matrix elements have been directly 
calculated at the physical quark masses, a big breakthrough compared with 
a few years ago. Still, the calculations were limited to  
the first couple leading moments. Higher moments, such as
$\langle x^2 \rangle$, have not been updated using dynamical fermions for more
than a decade~\cite{Gockeler:2004wp}. 
However, there are interesting proposals to obtain higher moments by using smeared sources to overcome the power-divergent mixing problem~\cite{Davoudi:2012ya} and by using light-quark--to--heavy-quark transition currents to compute current-current correlators in Euclidean space~\cite{Detmold:2005gg}.

Recently, one of the authors proposed a new approach to calculating the full 
$x$ dependence of parton quantities, such as the parton distributions and other
parton observables~\cite{Ji:2013dva}. The method is based on the observation
that, while in the rest frame of the nucleon, parton physics corresponds to
lightcone correlations, the same physics can be obtained through
time-independent spatial correlations in the infinite-momentum frame (IMF). For
finite but large momenta feasible in lattice simulations, a large-momentum
effective field theory (LaMET) can be used to relate Euclidean
quasi-distributions to physical ones through a factorization theorem~\cite{Ji:2014gla}.
Since then, there have been many follow-up calculations to determine the
one-loop corrections needed
to connect finite-momentum quasi-distributions to IMF/lightcone distributions 
for nonsinglet leading-twist PDFs~\cite{Xiong:2013bka}, generalized parton distributions (GPDs)~\cite{Ji:2015qla} and transversity GPDs~\cite{Xiong:2015nua} in the continuum. 
Reference~\cite{Ji:2015jwa} also explores the renormalization of quasi-distributions, and establishes that the quasi-distribution is multiplicatively renormalizable at two-loop order.  
There are also proposals to improve the quark correlators to remove linear divergences in one-loop matching~\cite{Li:2016amo} and to improve the nucleon source to get higher nucleon momenta on the lattice~\cite{Bali:2016lva}.
 
This new approach was immediately implemented on the lattice,
and first results of the technique were reported at various conferences in the summer of 2013~\cite{Lin:2014gaa,Lin:2014yra}.
Preliminary studies using the LaMET approach of the Bjorken-$x$ dependence of quark, helicity and transversity
distributions, along with the pion distribution amplitude, 
show reasonable signals for the quasi-distributions. 
In 2014, we reported the first attempt to make a lattice
calculation of the unpolarized isovector quark distributions using the LaMET
formalism~\cite{Lin:2014zya}. We use lattice gauge ensembles with $N_f=2+1+1$ highly improved staggered
quarks (HISQ) (generated by the MILC Collaboration) and
clover valence fermions with pion mass $310$~MeV. We establish the
convergence of the result within the
uncertainty of the calculation as the nucleon momentum increases. Although the lattice systematics are not yet
fully under control, we obtain some qualitative features of the flavor
structure of the nucleon sea: $\bar{d}>\bar{u}$.
In an independent follow-up lattice work, 
our result was confirmed by the ETMC Collaboration in 2015~\cite{Alexandrou:2015rja} using 
the twisted-mass fermion action. 
Reference~\cite{Lin:2014zya} also reports the total polarized sea asymmetry
$\Delta \bar{u}>\Delta \bar{d}$,
which was later confirmed in updated measurements by the STAR~\cite{Adamczyk:2014xyw} and PHENIX~\cite{Adare:2015gsd}
collaborations. 

In this work, we present the first lattice-QCD results for the helicity and transversity PDFs using the LaMET approach.
We will start by briefly reviewing the LaMET approach in Sec.~\ref{sec:LaMET}, 
and then discuss the finite-momentum corrections for quasi-distributions computed on the lattice in Sec.~\ref{sec:corrections}, using the unpolarized PDF as an example.
In Sec.~\ref{sec:PolPDFs}, we generalize the results to the spin-polarized PDFs, including helicity and transversity PDFs. Finally, we present the lattice results in Sec.~\ref{sec:num}, and discuss the
implications of these results~in Sec.~\ref{sec:discussion}, focusing on the
less-known antiquark distribution. The details of the finite-momentum corrections are given in the Appendices.
